%% ---------------------------------------------------------------------
%% 06 - KẾT LUẬN
%% ---------------------------------------------------------------------

\section{Kết luận}\label{sec:concl}

Nghiên cứu đã cài đặt thuật toán Funnel Tree cùng Thủ tục 2 (\emph{clip off
funnels}) bằng Python (NumPy cho mảng, không dùng thư viện hình học), kiểm chứng thuật toán trên bộ lưới đa
diện chuẩn, rồi đưa kết quả trắc địa vào một quy trình mô phỏng UAV trên mô hình
thành phố ba chiều.

Trên bộ lưới kiểm thử gồm 9 mesh (7 bề mặt kín với các dạng cầu và xuyến, 2 bề mặt
mở có cạnh biên), năm lưới có kết quả tham chiếu khớp tới bốn chữ số thập phân
trên tổng 76 đỉnh, sai số tương đối lớn nhất $1{,}7\times10^{-5}$; bốn mesh còn
lại chỉ được kiểm tra khả năng chạy và tính nhất quán nội tại. Nghiên cứu trường
hợp điển hình trên \texttt{cliff} cho thấy tuyến tới đỉnh vách cắt ngang một cạnh
tam giác tại điểm nội tại, đúng ưu thế của phương pháp trắc địa so với xấp xỉ qua
đỉnh; cơ chế cắt tỉa được bật trong mọi lần chạy và không làm sai kết quả trên bộ
dữ liệu có reference, nhưng repo không ghi log từng sự kiện clip nên bài không báo
số nút cây cụ thể.

Ở phần mô phỏng, ba bề mặt dẫn đường dựng từ mô hình thành phố \num{301490} tam
giác có \num{2295}, \num{3355} và \num{5913} đỉnh, tương ứng với ô danh nghĩa $5$,
$4$ và $3$ m. Trên lưới $4$ m, ba tuyến từ một nguồn đến ba đích dài
\num{142.909050}, \num{127.295935} và \num{86.194146} đơn vị mô phỏng. Chín
chuyến bay (ba lưới $\times$ ba đích) đều hợp lệ; bảy chuyến đạt mốc bám tuyến
$1$ m, trong đó lưới $4$ m là cấu hình duy nhất đạt cả ba với sai số lớn nhất
$0{,}399$--$0{,}601$ m. Lưới $5$ m và $3$ m mỗi lưới trượt đích G2 ($1{,}022$ m
và $1{,}381$ m): ca thứ nhất do waypoint không đều ở góc cua, ca thứ hai do vách
dốc $45^\circ$ đòi leo nhanh hơn giới hạn $5$ m/s $\cdot$ $3$ m/s$^2$. Kết quả
cho thấy việc chọn kích thước ô phải đi cùng dự phòng ngang, độ dốc và tham số
bay; không có cấu hình nào có sai số nhỏ nhất ở mọi đích.

Về vận hành, quy trình ba stage \texttt{data} $\to$ \texttt{plan} $\to$
\texttt{simulate} cùng chuỗi SHA và archive cho phép truy một con số về đúng scene,
bề mặt, mission, route và telemetry; image Gazebo Jetty được ghim digest nên kết
quả bay không đổi ngầm khi image đổi. Trong ba cấu hình, lưới $4$ m là lựa chọn
làm việc cho mission hiện tại; lưới $3$ m chỉ nên dùng sau khi giảm tốc độ hoặc
tăng gia tốc, vì vách $45^\circ$ vượt khả năng leo của profile $5$ m/s $\cdot$
$3$ m/s$^2$. Chi phí lập tuyến tăng nhanh hơn số tam giác ($45$ $\to$ $1\,008$
giây), nên lưới mịn hơn cần chạy ngoại tuyến hoặc tối ưu lan truyền.

Các kết quả xác nhận khả năng hoạt động của quy trình trong phạm vi thực nghiệm,
không tạo bảo đảm tối ưu cho đường bay trong không gian tự do ba chiều và không xác
nhận khả năng vận hành UAV thực tế. Các hướng phát triển tiếp theo gồm: lặp lại mỗi
cấu hình nhiều lần trong điều kiện tải được khống chế và giữ cố định điểm đầu cuối
khi so kích thước ô; thêm phép kiểm tra va chạm dọc spline và hiệu chỉnh tốc độ,
gia tốc cho vách dốc; tích hợp mô hình rotor hoặc autopilot cùng nhiễu trạng thái;
mở rộng sang nhiều nhiệm vụ và nhiều UAV có kiểm tra xung đột; và cải tiến hiệu
năng cài đặt Python cho lưới có hàng trăm nghìn tam giác.
